\documentclass[11pt,a4paper,oneside,openany]{book}
\usepackage{biblioflow-ista}

\renewcommand{\titolodocumento}{Pre-Modification Test Document}
\renewcommand{\sottotitolodocumento}{Pre-Modification Test Document}

% Prima colonna delle tabelle delle categorie: nome e descrizione.
\newcommand{\categoria}[3]{%
  \multirow{#3}{=}{\textbf{#1}\newline\footnotesize #2}}

\begin{document}
\frontmatter
\copertina
\tableofcontents
\listoftables

\mainmatter
\chapter{Introduzione}

\section{Scopo del Documento}

Lo scopo del presente documento è quello di presentare l'insieme dei casi di
test designati a realizzare il testing di sistema di BiblioFlow nella sua
versione attuale (Baseline), prima dell'introduzione delle modifiche previste
dalle Change Requests.

La strategia adottata è quella del testing black-box. Questo approccio
consente di verificare la conformità del sistema concentrandosi esclusivamente
sul comportamento esterno, ovvero sulla corrispondenza tra le richieste
inviate alle interfacce del sistema (API HTTP, middleware, provider di
accesso, funzioni delle automazioni) e le risposte e i dati prodotti, senza
ispezionare la logica interna del codice.

I casi descrivono il comportamento \emph{osservato} della versione di partenza,
non quello desiderato: dove il sistema presenta un difetto noto, l'oracolo
riporta il comportamento difettoso, così che la correzione introdotta dalle
Change Requests risulti visibile come differenza rispetto a questa base.

\section{Ambiente di Testing}

Al fine di garantire la coerenza e la riproducibilità dei risultati, i test di
sistema vengono eseguiti in un ambiente controllato configurato come segue:
\begin{itemize}
  \item \textbf{Versione del sistema:} tag \texttt{baseline-pre-cr-bf-01}
        (commit \texttt{e0742bf}). Poiché la baseline non contiene
        l'infrastruttura di test, l'esecuzione avviene sul commit
        \texttt{7469969}, che aggiunge Vitest e le fixture senza modificare i
        flussi verificati.
  \item \textbf{Runtime:} Node.js 22.12, npm 10.9, Vitest 4.1.11.
  \item \textbf{Database:} PostgreSQL 16 in un container effimero, database
        \texttt{biblioflow\_test}; la preparazione rifiuta qualunque indirizzo
        che non identifichi esplicitamente un database di test.
  \item \textbf{Interfaccia di input:} route handler e middleware invocati
        con richieste HTTP costruite nel test; funzioni delle automazioni
        invocate come farebbe l'attività pianificata.
  \item \textbf{Sistema Operativo di Riferimento:} Windows 11, fuso orario
        \texttt{Europe/Rome}.
\end{itemize}

\section{Tipologia di Testing}

Il metodo formale adottato per la progettazione dei casi di test è il
Category Partitioning (Partizione delle Categorie).

Questa tecnica è stata utilizzata per scomporre il dominio degli input del
sistema in unità funzionali distinte (categorie), identificando per ciascuna
di esse i parametri significativi e le relative classi di equivalenza. Tale
approccio ha permesso di coprire i flussi che le Change Requests andranno a
toccare (prenotazione, check-in, estensione, cancellazione, automazioni) e
quelli che dovranno restare invariati (accesso, area staff, notifiche),
minimizzando la ridondanza dei test.

\chapter{Identificazione delle Categorie e dei Parametri}

Il criterio adottato è quello del category partitioning, allo scopo di
individuare le categorie di input e i parametri a cui associare ciascuna di
esse. Di seguito sono elencati i gruppi di parametri definiti per il testing
di sistema di BiblioFlow, suddivisi per ambito logico: Accesso (chi effettua
la richiesta), Prenotazione (oggetto della richiesta) e Automazioni, Notifiche
e Realtime (elaborazioni periodiche e canali verso l'utente).

\section{Parametri di Accesso}

Fanno riferimento all'identità di chi invia la richiesta e alla risorsa a cui
tenta di accedere.

{\small
\begin{xltabular}{\linewidth}{L{5.4cm}L{1.3cm}Y}
  \toprule
  \intestazione \tb{Categoria} & \tb{Codice} & \tb{Scelte (Classi di Equivalenza)} \\
  \midrule
  \endhead
  \categoria{Sessione (SE)}{Presenza del cookie di sessione nella richiesta.}{2}
    & SE0 & Assente \\*
    & SE1 & Presente \\
  \midrule
  \categoria{Ruolo (RU)}{Ruolo dell'utente autenticato.}{3}
    & RU0 & Studente \\*
    & RU1 & Bibliotecario \\*
    & RU2 & Amministratore \\
  \midrule
  \categoria{Credenziali (CR)}{Esito del confronto della password all'accesso.}{2}
    & CR0 & Corrette \\*
    & CR1 & Password errata \\
  \midrule
  \categoria{Stato Account (AC)}{Abilitazione dell'account.}{2}
    & AC0 & Attivo \\*
    & AC1 & Disattivato \\
  \midrule
  \categoria{Risorsa (RS)}{Tipo di risorsa richiesta.}{2}
    & RS0 & Pagina dell'interfaccia \\*
    & RS1 & API \\
  \bottomrule
  \caption{Parametri relativi all'Accesso}\label{tab:par-accesso}\\
\end{xltabular}}

\section{Parametri della Prenotazione}

Fanno riferimento all'operazione richiesta sulla prenotazione e ai dati che la
accompagnano.

{\small
\begin{xltabular}{\linewidth}{L{5.4cm}L{1.3cm}Y}
  \toprule
  \intestazione \tb{Categoria} & \tb{Codice} & \tb{Scelte (Classi di Equivalenza)} \\
  \midrule
  \endhead
  \categoria{Operazione (OP)}{Operazione richiesta sulla prenotazione.}{6}
    & OP0 & Creazione \\*
    & OP1 & Cancellazione da parte dello studente \\*
    & OP2 & Check-in tramite l'endpoint dedicato \\*
    & OP3 & Check-in tramite \texttt{PATCH} (percorso legacy) \\*
    & OP4 & Estensione dell'orario \\*
    & OP5 & Annullamento da parte dello staff \\
  \midrule
  \categoria{Payload (PL)}{Completezza dei campi obbligatori.}{2}
    & PL0 & Completo \\*
    & PL1 & Campo obbligatorio mancante \\
  \midrule
  \categoria{Intervallo Orario (IO)}{Relazione tra l'intervallo richiesto e l'orario della sala.}{3}
    & IO0 & Ordinato e dentro l'orario di apertura \\*
    & IO1 & Fuori dall'orario di apertura \\*
    & IO2 & Invertito (fine precedente all'inizio) \\
  \midrule
  \categoria{Disponibilità (DI)}{Presenza di prenotazioni attive sovrapposte sullo stesso posto.}{2}
    & DI0 & Nessuna sovrapposizione \\*
    & DI1 & Sovrapposizione con una prenotazione attiva \\
  \midrule
  \categoria{Stato Iniziale (ST)}{Stato della prenotazione prima dell'operazione.}{3}
    & ST0 & \texttt{CONFERMATA} \\*
    & ST1 & \texttt{CHECK\_IN} \\*
    & ST2 & \texttt{SCADUTA} \\
  \midrule
  \categoria{Origine (OR)}{Chi produce la richiesta di prenotazione.}{2}
    & OR0 & Richiesta costruita direttamente \\*
    & OR1 & Payload prodotto dalla pagina \texttt{/prenota} \\
  \bottomrule
  \caption{Parametri relativi alla Prenotazione}\label{tab:par-prenotazione}\\
\end{xltabular}}

\section{Parametri di Automazioni, Notifiche e Realtime}

Fanno riferimento alle elaborazioni periodiche, alle funzioni dello staff e ai
canali con cui il sistema informa l'utente.

{\small
\begin{xltabular}{\linewidth}{L{5.4cm}L{1.3cm}Y}
  \toprule
  \intestazione \tb{Categoria} & \tb{Codice} & \tb{Scelte (Classi di Equivalenza)} \\
  \midrule
  \endhead
  \categoria{Momento (MO)}{Istante dell'elaborazione rispetto all'inizio della prenotazione.}{3}
    & MO0 & Dentro la finestra di check-in \\*
    & MO1 & Tra 15 e 20 minuti prima dell'inizio \\*
    & MO2 & Oltre il margine di no-show \\
  \midrule
  \categoria{Esecuzioni (ES)}{Numero di esecuzioni dell'automazione nella giornata.}{2}
    & ES0 & Una \\*
    & ES1 & Due \\
  \midrule
  \categoria{Funzione Staff (FS)}{Funzione dell'area staff invocata.}{4}
    & FS0 & Statistiche \\*
    & FS1 & Gestione utenti \\*
    & FS2 & Prestiti \\*
    & FS3 & Prenotazioni \\
  \midrule
  \categoria{Notifiche (NT)}{Operazione sulle notifiche.}{2}
    & NT0 & Avviso di posto liberato \\*
    & NT1 & Elenco filtrato e paginato \\
  \midrule
  \categoria{Canale SSE (CS)}{Fase della connessione Server-Sent Events.}{3}
    & CS0 & Apertura dello stream \\*
    & CS1 & Pubblicazione di un evento \\*
    & CS2 & Chiusura dello stream \\
  \bottomrule
  \caption{Parametri di Automazioni, Notifiche e Realtime}\label{tab:par-automazioni}\\
\end{xltabular}}

\chapter{Test di Sistema}

In questo capitolo vengono dettagliati i casi di test (Test Cases) progettati
per coprire i frame identificati. La tabella seguente mappa ogni Test Case
(ID) con il relativo Test Frame (combinazione di parametri) e l'Oracolo
(risultato atteso sulla versione di partenza). Gli oracoli segnati come
\emph{difetto} descrivono un comportamento errato della baseline, fissato di
proposito per misurare la correzione.

{\small
\begin{xltabular}{\linewidth}{L{2.4cm}L{3.4cm}Y}
  \toprule
  \intestazione \tb{TC ID} & \tb{Test Frame (Parametri)} & \tb{Oracolo / Risultato Atteso} \\
  \midrule
  \endhead
  TC-PRE-001 & SE1, RU0, OP0, PL0, IO0, DI0, OR0 &
  \textbf{Creazione valida:} HTTP 201, prenotazione \texttt{CONFERMATA},
  evento \texttt{PRENOTAZIONE\_CREATA} e notifica di conferma. \\
  TC-PRE-002 & SE1, RU0, OP0, PL1, OR0 &
  \textbf{Campo mancante:} HTTP 400 con l'elenco dei campi obbligatori;
  nessuna prenotazione creata. \\
  TC-PRE-003 & SE1, RU0, OP0, PL0, IO1, DI0 &
  \textbf{Fuori orario:} HTTP 400 con gli orari di apertura della sala;
  nessuna scrittura. \\
  TC-PRE-004 & SE1, RU0, OP0, PL0, IO2, DI0 &
  \textbf{Intervallo invertito (difetto):} HTTP 201; il sistema non verifica
  che la fine segua l'inizio. \\
  TC-PRE-005 & SE1, RU0, OP0, PL0, IO0, DI1 &
  \textbf{Posto occupato:} HTTP 409, \enquote{posto già prenotato}; nessuna
  nuova prenotazione. \\
  TC-PRE-006 & SE1, RU0, OP1, ST0 &
  \textbf{Cancellazione:} HTTP 200, stato \texttt{CANCELLATA}, evento
  \texttt{PRENOTAZIONE\_CANCELLATA}. \\
  TC-PRE-007 & SE1, RU0, OP2, ST0, MO0 &
  \textbf{Check-in valido rifiutato (difetto):} HTTP 400, \enquote{periodo di
  check-in scaduto}; il confronto usa un orario ancorato al 1970. \\
  TC-PRE-008 & SE1, RU0, OP3, ST0 &
  \textbf{Check-in legacy:} HTTP 200, prenotazione \texttt{CHECK\_IN}, posto
  \texttt{OCCUPATO}, evento \texttt{CHECK\_IN}. \\
  TC-PRE-009 & SE1, RU0, OP4, ST0, DI1 &
  \textbf{Estensione in conflitto:} HTTP 409; orario originale invariato. \\
  TC-PRE-010 & SE1, RU0, OP4, ST0, DI0 &
  \textbf{Estensione libera rifiutata (difetto):} HTTP 400 per durata
  superiore a otto ore, calcolata tra date incompatibili. \\
  TC-PRE-011 & CR0, AC0 &
  \textbf{Accesso valido:} utente restituito senza hash della password,
  ultimo accesso aggiornato. \\
  TC-PRE-012 & CR1, AC0 &
  \textbf{Password errata:} errore \enquote{Credenziali non valide}; ultimo
  accesso invariato. \\
  TC-PRE-013 & AC1 &
  \textbf{Account disattivato:} errore dedicato prima del confronto della
  password. \\
  TC-PRE-014 & SE0, RS0 &
  \textbf{Pagina staff anonima:} redirect 307 a \texttt{/login} con la
  destinazione in \texttt{callbackUrl}. \\
  TC-PRE-015 & SE0, RS1 &
  \textbf{API staff anonima:} HTTP 401 \enquote{Non autenticato}. \\
  TC-PRE-016 & SE1, RS0 &
  \textbf{Pagina con cookie:} il middleware lascia proseguire
  \texttt{/prenota} (verifica solo la presenza del cookie). \\
  TC-PRE-017 & SE1, RU0, RS1, FS0 &
  \textbf{Studente sulle statistiche:} HTTP 403, nessuna query di
  aggregazione. \\
  TC-PRE-018 & SE1, RU1, RS1, FS1 &
  \textbf{Bibliotecario disattiva un utente:} HTTP 403; operazione riservata
  all'amministratore, nessun aggiornamento. \\
  TC-PRE-019 & SE1, RU1, RS1, FS1 &
  \textbf{Dettaglio utente:} HTTP 200 con profilo e conteggi di
  prenotazioni, prestiti e no-show. \\
  TC-PRE-020 & SE1, RU2, RS1, FS2 &
  \textbf{Restituzione prestito:} stato \texttt{RESTITUITO}, data di
  restituzione, evento e notifica; zero giorni di ritardo. \\
  TC-PRE-021 & SE1, RU1, OP5, ST1, FS3 &
  \textbf{Annullamento staff:} prenotazione \texttt{CANCELLATA}, posto
  \texttt{DISPONIBILE}, evento e notifica. \\
  TC-PRE-022 & SE1, RU1, RS1, FS0 &
  \textbf{Tasso di no-show:} distribuzione 20\% no-show, 60\% completate,
  20\% altre. \\
  TC-PRE-023 & OP0, OR1 &
  \textbf{Contratto della pagina:} \texttt{POST /api/prenotazioni} con
  \texttt{userId}, posto, data, orari e margine pendolare. \\
  TC-PRE-024 & SE1, RU0, OP0, PL0, IO0, DI0, OR1 &
  \textbf{Conferma dalla pagina:} HTTP 201 e prenotazione
  \texttt{CONFERMATA}. \\
  PRE-AUT-001 & ST0, MO1, ES0 &
  \textbf{Promemoria:} una notifica \texttt{CHECK\_IN\_REMINDER} non letta
  con collegamento alla prenotazione ed evento \texttt{AUTOMATION}. \\
  PRE-AUT-002 & ST0, MO1, ES1 &
  \textbf{Promemoria idempotente:} la seconda esecuzione non invia nulla; un
  solo promemoria nel database. \\
  PRE-AUT-003 & ST0, MO2, ES0 &
  \textbf{No-show automatico:} stato \texttt{NO\_SHOW}, posto
  \texttt{DISPONIBILE}, notifica \texttt{ALERT}, evento
  \texttt{NO\_SHOW\_AUTO}. \\
  PRE-AUT-004 & ST2, MO2, ES0 &
  \textbf{Prenotazione già scaduta:} nessuna elaborazione, nessuna notifica
  né evento. \\
  PRE-NOT-001 & NT0 &
  \textbf{Posto liberato:} notifica \texttt{INFO} non letta, con collegamento
  a \texttt{/prenota}, all'utente che frequenta la sala. \\
  PRE-NOT-002 & SE1, RU0, NT1 &
  \textbf{Elenco notifiche:} HTTP 200 con \texttt{count=1},
  \texttt{totale=2}, \texttt{nonLette=1}. \\
  PRE-SSE-001 & CS0 &
  \textbf{Apertura stream:} HTTP 200, \texttt{text/event-stream}, cache
  disabilitata, commento iniziale \texttt{: connected}. \\
  PRE-SSE-002 & CS1 &
  \textbf{Isolamento dei canali:} l'evento sul canale dei posti arriva solo
  a quel canale e al canale jolly, nel formato \texttt{event}/\texttt{data}. \\
  PRE-SSE-003 & CS1 &
  \textbf{Aggiornamento posto:} evento \texttt{posto-update} con posto,
  stato, numero, sala e timestamp ISO. \\
  PRE-SSE-004 & CS2 &
  \textbf{Chiusura:} il client viene rimosso dal registro dell'emettitore. \\
  \bottomrule
\end{xltabular}}

Tutti i 34 casi sono stati eseguiti sulla versione di partenza con esito
conforme all'oracolo (34 su 34). I tre casi segnati come difetto (TC-PRE-004,
TC-PRE-007, TC-PRE-010) e il contratto di TC-PRE-023 costituiscono l'evidenza
da cui partono le Change Requests.

\end{document}
